\clearpage
\section{Rootwise all vector connection
estimates}\label{e-rootwise-all-vector-connection-estimates}

Technical derivation for review. Equation and subsection numbers in this
appendix are local to the source derivation. Historical local labels are
retained, including numbering gaps or nonsequential labels. The
accompanying source notes retain the original expressions for
comparison.

\subsection{1. Assumptions and the excitation
energy}\label{e-assumptions-and-the-excitation-energy}

Fix the rank and a partition. For every root/pair \(e=(a,b)\), let
\(r_{e}=\vert z_{a}-z_{b}\vert\) and assume
\(\alpha _{a}+\alpha _{b} \le  \kappa  r_{e}\), with \(\kappa\)
sufficiently small. Write \(r_{*}=\min _{e} r_{e}\) and
\(W_{2}=\sum _{e} r_{e}^{-2}\). No ratio \(r_{\max}/r_{*}\) is
restricted.

Freeze the slow variables. Let

% DISPLAY e.1
% SOURCE H_fast,e = H_B,e + H_F,e,
% SOURCE E_0,e = E_B,e + E_F,min,e,
% SOURCE h_e = H_fast,e - E_0,e
% SOURCE     = (H_B,e-E_B,e) + (H_F,e-E_F,min,e).
\begin{gather*}
\begin{aligned}
& H_{\mathrm{fast},e} = H_{B,e} + H_{F,e}, \\
& E_{0,e} = E_{B,e} + E_{F,\mathrm{min},e}, \\
& h_{e} = H_{\mathrm{fast},e} - E_{0,e} \\
& = (H_{B,e}-E_{B,e}) + (H_{F,e}-E_{F,\mathrm{min},e}).
\end{aligned}\notag
\end{gather*}

Here \(H_{B,e}\) is the exact kinetic-aware bosonic oscillator. The
uniform pair-matrix bounds give

% DISPLAY e.2
% SOURCE c(p_e^2+r_e^2 Y_e^2) <= H_B,e <= C(p_e^2+r_e^2 Y_e^2),
% SOURCE E_B,e <= C r_e,
% SOURCE h_e >= c r_e Q_e,
\begin{gather*}
\begin{aligned}
& c(p_{e}^{2}+r_{e}^{2} Y_{e}^{2}) \le  H_{B,e} \le  C(p_{e}^{2}+r_{e}^{2} Y_{e}^{2}), \\
& E_{B,e} \le  C r_{e}, \\
& h_{e} \ge  c r_{e} Q_{e},
\end{aligned}\notag
\end{gather*}

where \(P_{e}\) is its full boson/fermion vacuum projection and
\(Q_{e}=1-P_{e}\). In particular

% DISPLAY e.3
% SOURCE ||p_e Q_e v||^2 <= C h_e[v],
% SOURCE ||Y_e v||^2 <= C r_e^-2 h_e[v] + C r_e^-1 ||v||^2,
% SOURCE ||p_e P_e v||^2 <= C r_e ||P_e v||^2,
% SOURCE ||Y_e p_e P_e v||^2 <= C ||P_e v||^2.                 (1)
\begin{gather*}
\begin{aligned}
& \Vert p_{e} Q_{e} v\Vert ^{2} \le  C h_{e}[v], \\
& \Vert Y_{e} v\Vert ^{2} \le  C r_{e}^{-2} h_{e}[v] + C r_{e}^{-1} \Vert v\Vert ^{2}, \\
& \Vert p_{e} P_{e} v\Vert ^{2} \le  C r_{e} \Vert P_{e} v\Vert ^{2}, \\
& \Vert Y_{e} p_{e} P_{e} v\Vert ^{2} \le  C \Vert P_{e} v\Vert ^{2}.
\end{aligned}\tag{1}
\end{gather*}

The last two statements also hold with any uniformly bounded finite
color/spatial matrices between the displayed operators. Constants depend
only on rank and the fixed small pair-matrix neighborhood. \(h_{e}\),
rather than \(H_{B,e}\) or a centered oscillator, is essential.

Let \(\chi\) be any multiplication cutoff with
\(\vert \chi \vert \le 1\) supported in
\(\vert Y\vert \le \sigma  r_{*}\). Derivatives of \(\chi\) are not used
below. Every estimate is a fast-fiber estimate for arbitrary vectors,
with arbitrary spectator fast modes and arbitrary slow/Clifford
coefficient vectors. It can be integrated over the slow variables. No
internal momentum or oscillator-degree cutoff is imposed. No individual
internal singlet condition is imposed; the estimates remain valid on the
total residual-equivariant space.

\subsection{2. Same-root and mixed-root
lemmas}\label{e-same-root-and-mixed-root-lemmas}

For any pair \(e\) and any positive denominator scale \(r_{g}\) drawn
from the retained roots,

% DISPLAY e.4
% SOURCE ||chi Y_e p_e v/r_g||^2
% SOURCE   <= C sigma^2 (r_*^2/r_g^2) h_e[v]
% SOURCE        + C r_g^-2 ||v||^2.                           (2)
\begin{gather*}
\begin{aligned}
& \Vert \chi  Y_{e} p_{e} v/r_{g}\Vert ^{2} \\
& \le  C \sigma ^{2} (r_{*}^{2}/r_{g}^{2}) h_{e}[v] \\
& + C r_{g}^{-2} \Vert v\Vert ^{2}.
\end{aligned}\tag{2}
\end{gather*}

Proof: split the INPUT as \(P_{e} v+Q_{e} v\). On \(Q_{e}\) use
\(\vert \chi  Y_{e}\vert \le \sigma  r_{*}\) and the first inequality in
(1). On \(P_{e}\) use the last inequality in (1), dropping \(\chi\). The
factor two from the split is harmless. \(P_{e} v\) need not be supported
in the tube. No argument asserts that \(P_{e}\) preserves tube support.

For distinct roots \(e\) and \(f\),

% DISPLAY e.5
% SOURCE ||chi Y_f p_e v/r_g||^2
% SOURCE   <= C sigma^2 (r_*^2/r_g^2) h_e[v]
% SOURCE    + C r_e/(r_g^2 r_f^2) h_f[v]
% SOURCE    + C r_e/(r_g^2 r_f) ||v||^2.                      (3)
\begin{gather*}
\begin{aligned}
& \Vert \chi  Y_{f} p_{e} v/r_{g}\Vert ^{2} \\
& \le  C \sigma ^{2} (r_{*}^{2}/r_{g}^{2}) h_{e}[v] \\
& + C r_{e}/(r_{g}^{2} r_{f}^{2}) h_{f}[v] \\
& + C r_{e}/(r_{g}^{2} r_{f}) \Vert v\Vert ^{2}.
\end{aligned}\tag{3}
\end{gather*}

Indeed the \(Q_{e}\) input is as above. On the \(P_{e}\) input, drop
\(\chi\), commute \(Y_{f}\) with \(P_{e}\), and use
\(\Vert p_{e} P_{e}\Vert ^{2}\le C r_{e}\) followed by the \(Y_{f}\)
bound in (1). Equivalently split the remaining \(f\) input into
\(P_{f}\) and \(Q_{f}\). This is the step that cannot be omitted from a
mixed-root claim.

If e,f,g are the three edges of a triangle, \(r_{e}\le r_{f}+r_{g}\), so

% DISPLAY e.6
% SOURCE r_e/(r_g^2 r_f^2) <= 2 r_*^-3,
% SOURCE r_e/(r_g^2 r_f) <= r_g^-2+(r_g r_f)^-1 <= C W_2,t.
\begin{gather*}
\begin{aligned}
& r_{e}/(r_{g}^{2} r_{f}^{2}) \le  2 r_{*}^{-3}, \\
& r_{e}/(r_{g}^{2} r_{f}) \le  r_{g}^{-2}+(r_{g} r_{f})^{-1} \le  C W_{2,t}.
\end{aligned}\notag
\end{gather*}

Consequently

% DISPLAY e.7
% SOURCE ||chi Y_f p_e v/r_g||^2
% SOURCE   <= C (sigma^2+r_*^-3) (h_e[v]+h_f[v])
% SOURCE        + C W_2,t ||v||^2.                            (4)
\begin{gather*}
\begin{aligned}
& \Vert \chi  Y_{f} p_{e} v/r_{g}\Vert ^{2} \\
& \le  C (\sigma ^{2}+r_{*}^{-3}) (h_{e}[v]+h_{f}[v]) \\
& + C W_{2,t} \Vert v\Vert ^{2}.
\end{aligned}\tag{4}
\end{gather*}

For the input-root denominator \(g=e\), the same conclusion follows
directly, without any triangle hypothesis, with residual
\((r_{e} r_{f})^{-1}\). In (2)-(4), \(Y\) and \(p\) can be contracted
with any bounded root-preserving finite matrices. Finite sums of such
sources obey the same bounds after enlarging the rank-dependent
constant.

A triangle hypothesis (or another restriction on \(r_{e}\) relative to
\(r_{f},r_{g}\)) is genuinely necessary to replace (3) by (4): the
\(P_{e} P_{f}\) input has squared size comparable to
\(r_{e}/(r_{g}^{2} r_{f})\). An arbitrary unrelated \(r_{e}\) cannot
disappear from this expression.

\subsection{3. Exact inverse-M factorization for the fast
column}\label{e-exact-inverse-m-factorization-for-the-fast-column}

Use the notation of the intrinsic geometry and adapted source inventory:

% DISPLAY e.8
% SOURCE M=F+P,   P=L-J_2,   J_2=K C_Y=R^-1 C_Y* C_Y,
% SOURCE W=I+K K*,   K=R^-1 C_Y*,
% SOURCE T=(I+F^-1 P*)^-1,
% SOURCE M^-*=T F^-1,
% SOURCE d_0(v)=F^-1 C* p_F v.
\begin{gather*}
\begin{aligned}
& M=F+P, \quad  P=L-J_{2}, \quad  J_{2}=K C_{Y}=R^{-1} C_{Y}^{*} C_{Y}, \\
& W=I+K K^{*}, \quad  K=R^{-1} C_{Y}^{*}, \\
& T=(I+F^{-1} P^{*})^{-1}, \\
& M^{-*}=T F^{-1}, \\
& d_{0}(v)=F^{-1} C^{*} p_{F} v.
\end{aligned}\notag
\end{gather*}

\(F\) is self-adjoint pair-block diagonal;
\(\Vert F_{e}^{-1} C_{e}^{*}\Vert \le C \kappa\). On the tube,

% DISPLAY e.9
% SOURCE ||P F^-1||+||F^-1 P*|| <= C sigma,
% SOURCE ||T||+||W^(1/2)|| <= C,
% SOURCE ||K|| <= C sigma.
\begin{gather*}
\begin{aligned}
& \Vert P F^{-1}\Vert +\Vert F^{-1} P^{*}\Vert  \le  C \sigma , \\
& \Vert T\Vert +\Vert W^{1/2}\Vert  \le  C, \\
& \Vert K\Vert  \le  C \sigma .
\end{aligned}\notag
\end{gather*}

Let \(B p_{F}=W^{1/2} M^{-*}(C^{*}+N_{Y}^{*})p_{F}\) be the exact fast
column of the Gauss stack. Algebraically, with the indicated operator
ordering,

% DISPLAY e.10
% SOURCE B p_F-d_0
% SOURCE   =(W^(1/2)-I)d_0
% SOURCE    +W^(1/2)T[
% SOURCE       F^-1 N_Y* p_F
% SOURCE       -F^-1 L* d_0
% SOURCE       +F^-1 J_2* d_0].                               (5)
\begin{gather*}
\begin{aligned}
& B p_{F}-d_{0} \\
& =(W^{1/2}-I)d_{0} \\
& +W^{1/2}T[ \\
& F^{-1} N_{Y}^{*} p_{F} \\
& -F^{-1} L^{*} d_{0} \\
& +F^{-1} J_{2}^{*} d_{0}].
\end{aligned}\tag{5}
\end{gather*}

There are no derivatives of inverse matrices in this first-order form
identity.

The first two bracketed expressions are finite triangle sources
\(Y_{f} p_{e}/r_{g}\), with respectively bounded and \(O(\kappa )\)
root-preserving factors. The last expression factors as

% DISPLAY e.11
% SOURCE F^-1 J_2* d_0
% SOURCE   =F^-1 C_Y* [C_Y R^-1 d_0],                         (6)
\begin{gather*}
\begin{aligned}
& F^{-1} J_{2}^{*} d_{0} \\
& =F^{-1} C_{Y}^{*} [C_{Y} R^{-1} d_{0}],
\end{aligned}\tag{6}
\end{gather*}

where \(\Vert F^{-1} C_{Y}^{*}\Vert \le C \sigma\) and
\(C_{Y} R^{-1} d_{0}\) is a sum of same-root terms \(Y_{e} p_{e}/r_{e}\)
with \(O(\kappa )\) coefficients. Also

% DISPLAY e.12
% SOURCE (W^(1/2)-I)d_0
% SOURCE   =(W^(1/2)+I)^-1 K [K* d_0],
% SOURCE K* d_0=C_Y R^-1 d_0.                                 (7)
\begin{gather*}
\begin{aligned}
& (W^{1/2}-I)d_{0} \\
& =(W^{1/2}+I)^{-1} K [K^{*} d_{0}], \\
& K^{*} d_{0}=C_{Y} R^{-1} d_{0}.
\end{aligned}\tag{7}
\end{gather*}

Thus all potentially troublesome inverse-M matrices remain bounded left
factors. Equations (2)-(7) prove

% DISPLAY e.13
% SOURCE ||chi (B p_F-d_0)v||^2
% SOURCE   <= C (sigma^2+r_*^-3) sum_e h_e[v]
% SOURCE        + C W_2 ||v||^2.                              (8)
\begin{gather*}
\begin{aligned}
& \Vert \chi  (B p_{F}-d_{0})v\Vert ^{2} \\
& \le  C (\sigma ^{2}+r_{*}^{-3}) \sum _{e} h_{e}[v] \\
& + C W_{2} \Vert v\Vert ^{2}.
\end{aligned}\tag{8}
\end{gather*}

The estimate holds for arbitrary fast vectors. It is not an estimate
restricted to vacuum or a finite excited coefficient space.

The spin column obeys separately

% DISPLAY e.14
% SOURCE ||chi W^(1/2)M^-* S_f v||^2 <= C W_2 ||v||^2,          (9)
\begin{gather*}
\begin{aligned}
& \Vert \chi  W^{1/2}M^{-*} S_{f} v\Vert ^{2} \le  C W_{2} \Vert v\Vert ^{2},
\end{aligned}\tag{9}
\end{gather*}

because the finite fermionic Gauss matrices are bounded and \(F^{-1}\)
retains the individual output-root denominators.

\subsection{4. Exact slow-square coefficient and its Schur
loss}\label{e-exact-slow-square-coefficient-and-its-schur-loss}

Let \(G_{s}=\operatorname{diag} (g_{C}^{-1},I_{I})\), so the direct slow
kinetic stack is \(w=G_{s}^{1/2}p_{s}\). The exact normalized slow
column in the Gauss square is

% DISPLAY e.15
% SOURCE A=W^(1/2)M^-* Psi,
% SOURCE Psi=[(C_Y*+(F+L)*K)g_C^-1/2, I_Y*].                  (10)
\begin{gather*}
\begin{aligned}
& A=W^{1/2}M^{-*} \Psi , \\
& \Psi =[(C_{Y}^{*}+(F+L)^{*}K)g_{C}^{-1/2}, I_{Y}^{*}].
\end{aligned}\tag{10}
\end{gather*}

This follows directly from \(O^{*}Sg_{S}^{-1}\). In particular the
\(F K\) term must be retained. One has

% DISPLAY e.16
% SOURCE ||Psi*F^-1|| <= C sigma,
% SOURCE ||A|| <= C sigma.                                    (11)
\begin{gather*}
\begin{aligned}
& \Vert \Psi ^{*}F^{-1}\Vert  \le  C \sigma , \\
& \Vert A\Vert  \le  C \sigma .
\end{aligned}\tag{11}
\end{gather*}

The adjoint coefficient acting on the frozen fast column has the exact
expression

% DISPLAY e.17
% SOURCE A* d_0=Psi*F^-1 T*W^(1/2)d_0.                        (12)
\begin{gather*}
\begin{aligned}
& A^{*} d_{0}=\Psi ^{*}F^{-1} T^{*}W^{1/2}d_{0}.
\end{aligned}\tag{12}
\end{gather*}

Expand

% DISPLAY e.18
% SOURCE T*W^(1/2)-I=(T*-I)+T*(W^(1/2)-I),
% SOURCE T*-I=-T* P F^-1.                                    (13)
\begin{gather*}
\begin{aligned}
& T^{*}W^{1/2}-I=(T^{*}-I)+T^{*}(W^{1/2}-I), \\
& T^{*}-I=-T^{*} P F^{-1}.
\end{aligned}\tag{13}
\end{gather*}

The leading term \(\Psi ^{*}F^{-1} d_{0}\) is the stack of

% DISPLAY e.19
% SOURCE I_Y F^-1 d_0,
% SOURCE g_C^-1/2[C_Y F^-1 d_0+K* d_0+K* L F^-1 d_0].         (14)
\begin{gather*}
\begin{aligned}
& I_{Y} F^{-1} d_{0}, \\
& g_{C}^{-1/2}[C_{Y} F^{-1} d_{0}+K^{*} d_{0}+K^{*} L F^{-1} d_{0}].
\end{aligned}\tag{14}
\end{gather*}

The first three contractions are same-root \(Y_{e} p_{e}/r_{e}\).
\(L F^{-1} d_{0}\) is a triangle source with INPUT-root denominator
\(r_{e}\), so (3) applies even without using a frequency ratio. The
\(J_{2}\) part of \(P F^{-1} d_{0}\) factors as

% DISPLAY e.20
% SOURCE J_2 F^-1 d_0=R^-1 C_Y* [C_Y F^-1 d_0],               (15)
\begin{gather*}
\begin{aligned}
& J_{2} F^{-1} d_{0}=R^{-1} C_{Y}^{*} [C_{Y} F^{-1} d_{0}],
\end{aligned}\tag{15}
\end{gather*}

again a bounded \(O(\sigma )\) left factor times a same-root
contraction. Equations (7), (11), and (13)-(15) therefore prove the
sharper estimate

% DISPLAY e.21
% SOURCE ||chi A* d_0(v)||^2
% SOURCE   <= C kappa^2 (sigma^2+r_*^-3) sum_e h_e[v]
% SOURCE        + C kappa^2 W_2 ||v||^2.                      (16)
\begin{gather*}
\begin{aligned}
& \Vert \chi  A^{*} d_{0}(v)\Vert ^{2} \\
& \le  C \kappa ^{2} (\sigma ^{2}+r_{*}^{-3}) \sum _{e} h_{e}[v] \\
& + C \kappa ^{2} W_{2} \Vert v\Vert ^{2}.
\end{aligned}\tag{16}
\end{gather*}

Combining (8), (9), and \(\Vert A\Vert \le C \sigma\) also controls
\(A^{*}\) applied to the full fast-plus-spin stack, with a coefficient
\(C(\kappa ^{2}+\sigma ^{2})\) multiplying the right side of (8), and a
harmless \(C \sigma ^{2} W_{2}\) spin term.

For completeness, for any vectors w,b and \(\eta >0\),

% DISPLAY e.22
% SOURCE eta||w||^2+||b-Aw||^2
% SOURCE   =||(eta I+A*A)^(1/2)
% SOURCE       [w-(eta I+A*A)^-1 A*b]||^2
% SOURCE       +<b,(I+eta^-1 A A*)^-1 b>.
\begin{gather*}
\begin{aligned}
& \eta \Vert w\Vert ^{2}+\Vert b-Aw\Vert ^{2} \\
& =\Vert (\eta  I+A^{*}A)^{1/2} \\
& [w-(\eta  I+A^{*}A)^{-1} A^{*}b]\Vert ^{2} \\
& +\langle b,(I+\eta ^{-1} A A^{*})^{-1} b\rangle .
\end{aligned}\notag
\end{gather*}

The difference between \(\Vert b\Vert ^{2}\) and the final term is at
most \(\eta ^{-1}\Vert A^{*}b\Vert ^{2}\). Thus (16) controls the
frozen-column Schur loss by exact excitation energies and \(W_{2}\),
with no \(\delta  \sum _{e} r_{e}\) zero-point cost. This is a
first-order form identity: variable coefficients create no omitted
ordering terms at this stage.

\subsection{5. What this lemma does and does not
close}\label{e-what-this-lemma-does-and-does-not-close}

It closes the all-vector rootwise CONNECTION-NORM sublemma, including
mixed roots, exact inverse-M left factors, and the coefficient occurring
after slow-square completion. It remains valid with unrestricted
internal momenta because none of its steps estimates those momenta or
differentiates a fast projection in a slow direction.

At the connection-norm stage, two additional arguments are necessary.
Section \(6\) below supplies the first for the fast kinetic column:

\begin{enumerate}
\def\labelenumi{\arabic{enumi}.}
\tightlist
\item
  Bounding a squared connection norm is not a bound for the bilinear
  pairing with \(d_{0}\). The direct Young estimate against
  \(\Vert d_{0} v\Vert ^{2}\) introduces its frozen vacuum energy
  \(C \kappa ^{2} \sum _{e} r_{e}\). The odd H1 source / normal-ordering
  or graph-local form argument must be used instead; (8) alone does not
  perform it.
\item
  Completing all slow squares leaves shifted slow derivatives. Splitting
  off \(\eta\) times ordinary internal kinetic energy is not the same as
  splitting \(\eta\) times the colored internal supersymmetric form: an
  unweighted Yukawa remainder \(\eta  C \alpha _{a}\) may be large for a
  far block. One must retain a suitable full slow-supercharge form or
  prove an incident-weighted conversion of the shifted slow squares.
  Equations (10)-(16) do not silently supply that conversion.
\end{enumerate}

Accordingly these results are compatible with, but do not establish, a
near-sharp complementary lower bound in the original Hamiltonian. They
provide the rootwise estimate that avoids a global-frequency
substitution at the connection stage.

\subsection{6. Addendum: signed fast kinetic terms without a
vacuum-frequency
loss}\label{e-addendum-signed-fast-kinetic-terms-without-a-vacuum-frequency-loss}

The first obstruction listed in Section \(5\) can also be removed for
the exact fast kinetic column. The following statements concern a smooth
core vector \(v\) supported in the fast tube; an outer cutoff equal to
one near its support may always be used in Sections 2-4. Only that
support requirement, not a restriction on the oscillator content of
\(v\), is imposed.

\subsubsection{6.1 All-vector triangle
multiplication}\label{e-all-vector-triangle-multiplication}

For distinct triangle roots e,f,g, write \(s_{t}=r_{e}+r_{f}+r_{g}\).
Any uniformly bounded real coefficient tensor that may depend on the
slow variables but is independent of the fast variables \(Y\) obeys

% DISPLAY e.23
% SOURCE |<v,s_t Y_e Y_f Y_g v>|
% SOURCE    <= C sigma sum_(j in t) h_j[v]
% SOURCE       + C sigma^-1 W_2,t ||v||^2.                   (17)
\begin{gather*}
\begin{aligned}
& \vert \langle v,s_{t} Y_{e} Y_{f} Y_{g} v\rangle \vert \\
& \le  C \sigma  \sum _{j \in  t} h_{j}[v] \\
& + C \sigma ^{-1} W_{2,t} \Vert v\Vert ^{2}.
\end{aligned}\tag{17}
\end{gather*}

First, splitting inputs into \(Q_{f}, P_{f} Q_{g}\), and \(P_{f} P_{g}\)
gives

% DISPLAY e.24
% SOURCE ||chi Y_f Y_g v||^2
% SOURCE    <= C sigma^2 r_*^2 [r_f^-2 h_f[v]+r_g^-2 h_g[v]]
% SOURCE       + C/(r_f r_g) ||v||^2.                        (18)
\begin{gather*}
\begin{aligned}
& \Vert \chi  Y_{f} Y_{g} v\Vert ^{2} \\
& \le  C \sigma ^{2} r_{*}^{2} [r_{f}^{-2} h_{f}[v]+r_{g}^{-2} h_{g}[v]] \\
& + C/(r_{f} r_{g}) \Vert v\Vert ^{2}.
\end{aligned}\tag{18}
\end{gather*}

On the first component bound \(\chi  Y_{g}\) by \(\sigma  r_{*}\) and
use the excited \(Y_{f}\) oscillator bound; on the second use
\(\chi  Y_{f}\) and the excited \(Y_{g}\) bound; on the last use both
Gaussian widths. The splitting is for estimation only; the projected
pieces need not have tube support.

Choose \(e\) to be a largest triangle root. If
\(U_{B,e}=\text{ constant } \exp (-Y_{e} \cdot  \Omega _{e} Y_{e}/2)\),
put \(a_{e}=\partial _{e}+\Omega _{e} Y_{e}\). The exact bosonic
excitation form controls \(\Vert a_{e} v\Vert ^{2}\). Since the other
two factors do not depend on \(Y_{e}\), integration by parts gives,
componentwise,

% DISPLAY e.25
% SOURCE <v,Y_e Y_f Y_g v>
% SOURCE    =Re<a_e v,Omega_e^-1(Y_f Y_g)v>.
\begin{gather*}
\begin{aligned}
& \langle v,Y_{e} Y_{f} Y_{g} v\rangle \\
& =\operatorname{Re} \langle a_{e} v,\Omega _{e}^{-1}(Y_{f} Y_{g})v\rangle .
\end{aligned}\notag
\end{gather*}

Using \(\Vert \Omega _{e}^{-1}\Vert \le C/r_{e}, s_{t}/r_{e}\le 3\), and
(18) gives (17) by Young with parameter \(\sigma\). The proof is valid
for arbitrary fermionic and slow vector coefficients.

\subsubsection{6.2 Leading triangle differential
forms}\label{e-leading-triangle-differential-forms}

Let a(Y) be a real symmetric fast coefficient matrix supported only in
the off-diagonal root blocks \((e,g)\) and \((g,e)\), with each entry a
bounded linear contraction of \(Y_{f}\) divided by \(r_{e}\), multiplied
by a coefficient \(O(\kappa )\) that is independent of the fast
variables \(Y\) (it may depend on the slow variables). For the product
bosonic Gaussian \(U_{B}\), writing \(v=U_{B} u\) gives the exact
identity

% DISPLAY e.26
% SOURCE integral (partial v)* a (partial v)
% SOURCE   =integral U_B^2 (partial u)* a (partial u)
% SOURCE    +integral [div(a Omega Y)-(Omega Y)·a(Omega Y)] |v|^2.
% SOURCE                                                                 (19)
\begin{gather*}
\begin{aligned}
& \int  (\partial  v)^{*} a (\partial  v) \\
& =\int  U_{B}^{2} (\partial  u)^{*} a (\partial  u) \\
& +\int  [\operatorname{div} (a \Omega  Y)-(\Omega  Y) \cdot  a(\Omega  Y)] \vert v\vert ^{2}.
\end{aligned}\tag{19}
\end{gather*}

This identity does not project onto the fermionic vacuum. The excitation
form satisfies

% DISPLAY e.27
% SOURCE sum_e (H_B,e-E_B,e)[v]
% SOURCE    =integral U_B^2 (partial u)*G(partial u),
\begin{gather*}
\begin{aligned}
& \sum _{e} (H_{B,e}-E_{B,e})[v] \\
& =\int  U_{B}^{2} (\partial  u)^{*}G(\partial  u),
\end{aligned}\notag
\end{gather*}

with \(G\) uniformly positive. On tube-supported \(v\), the first term
in (19) is bounded in absolute value by
\(C \kappa  \sigma  \sum _{e} h_{e}[v]\). The divergence term in (19)
vanishes: a connects distinct roots, depends on the third root, and
\(\Omega\) is pair-block diagonal. The remaining scalar is a triangle
cubic with coefficient \(O(\kappa  r_{g})\), hence is bounded by (17).
Consequently

% DISPLAY e.28
% SOURCE |integral (partial v)* a (partial v)|
% SOURCE   <= C kappa sigma sum_(j in t) h_j[v]
% SOURCE       + C kappa sigma^-1 W_2,t ||v||^2.              (20)
\begin{gather*}
\begin{aligned}
& \vert \int  (\partial  v)^{*} a (\partial  v)\vert \\
& \le  C \kappa  \sigma  \sum _{j \in  t} h_{j}[v] \\
& + C \kappa  \sigma ^{-1} W_{2,t} \Vert v\Vert ^{2}.
\end{aligned}\tag{20}
\end{gather*}

The same proof applies to \(O(\kappa ^{2})\) coefficients. Equation (20)
applies to both leading triangle forms from

% DISPLAY e.29
% SOURCE 2 Re<d_0,F^-1 N_Y* p_F-F^-1 L* d_0>.
\begin{gather*}
\begin{aligned}
& 2 \operatorname{Re} \langle d_{0},F^{-1} N_{Y}^{*} p_{F}-F^{-1} L^{*} d_{0}\rangle .
\end{aligned}\notag
\end{gather*}

No derivative cutoff is needed: (19) is applied to the original
polynomial coefficient and \(v\) itself has tube support. If one instead
wants a globally cut-off form on arbitrary unsupported inputs, cutoff
derivatives must be included; an even radial cutoff generates additional
triangle cubics, with coefficient smaller by
\(O((\sigma ^{2} r_{*}^{3})^{-1})\).

\subsubsection{6.3 Full inverse-M
reduction}\label{e-full-inverse-m-reduction}

There is a short exact reduction avoiding an infinite expansion. Set

% DISPLAY e.30
% SOURCE Z=F^-1 P*,   T=(I+Z)^-1,
% SOURCE n=F^-1 N_Y* p_F,
% SOURCE X=n-Z d_0.
\begin{gather*}
\begin{aligned}
& Z=F^{-1} P^{*}, \quad  T=(I+Z)^{-1}, \\
& n=F^{-1} N_{Y}^{*} p_{F}, \\
& X=n-Z d_{0}.
\end{aligned}\notag
\end{gather*}

Then

% DISPLAY e.31
% SOURCE B p_F=W^(1/2)(d_0+T X).
\begin{gather*}
\begin{aligned}
& B p_{F}=W^{1/2}(d_{0}+T X).
\end{aligned}\notag
\end{gather*}

Because \(W\ge I\) and \(T-I=-ZT\),

% DISPLAY e.32
% SOURCE ||B p_F v||^2
% SOURCE   >=||d_0 v||^2+2 Re<d_0 v,Xv>
% SOURCE       -2 |<Z* d_0 v,T Xv>|.                         (21)
\begin{gather*}
\begin{aligned}
& \Vert B p_{F} v\Vert ^{2} \\
& \ge \Vert d_{0} v\Vert ^{2}+2 \operatorname{Re} \langle d_{0} v,Xv\rangle \\
& -2 \vert \langle Z^{*} d_{0} v,T Xv\rangle \vert .
\end{aligned}\tag{21}
\end{gather*}

Both \(X\) and \(Z^{*} d_{0}=P F^{-1} d_{0}\) are covered by the
rootwise connection estimates. The last pairing is therefore bounded by

% DISPLAY e.33
% SOURCE C(sigma^2+r_*^-3) sum_e h_e[v]+C W_2||v||^2.
\begin{gather*}
\begin{aligned}
& C(\sigma ^{2}+r_{*}^{-3}) \sum _{e} h_{e}[v]+C W_{2}\Vert v\Vert ^{2}.
\end{aligned}\notag
\end{gather*}

Inside \(\langle d_{0},X\rangle\), the \(J_{2}\) term factors exactly as

% DISPLAY e.34
% SOURCE <d_0,F^-1 J_2* d_0>
% SOURCE    =<C_Y F^-1 d_0,C_Y R^-1 d_0>,                    (22)
\begin{gather*}
\begin{aligned}
& \langle d_{0},F^{-1} J_{2}^{*} d_{0}\rangle \\
& =\langle C_{Y} F^{-1} d_{0},C_{Y} R^{-1} d_{0}\rangle ,
\end{aligned}\tag{22}
\end{gather*}

and has the same bound. Only the two linear triangle forms in Section
\(6.2\) remain. Combining (20)-(22) gives

% DISPLAY e.35
% SOURCE ||B p_F v||^2
% SOURCE   >=||d_0 v||^2
% SOURCE      -C[kappa sigma+sigma^2+r_*^-3] sum_e h_e[v]
% SOURCE      -C(1+kappa/sigma) W_2||v||^2.                  (23)
\begin{gather*}
\begin{aligned}
& \Vert B p_{F} v\Vert ^{2} \\
& \ge \Vert d_{0} v\Vert ^{2} \\
& -C[\kappa  \sigma +\sigma ^{2}+r_{*}^{-3}] \sum _{e} h_{e}[v] \\
& -C(1+\kappa /\sigma ) W_{2}\Vert v\Vert ^{2}.
\end{aligned}\tag{23}
\end{gather*}

Thus direct \(p_{F}^{2}\) plus the exact Gauss FAST column, the adapted
quadratic potential, and its fast fermion mass have a lower bound

% DISPLAY e.36
% SOURCE E_0,total ||v||^2
% SOURCE    +(1-C[kappa sigma+sigma^2+r_*^-3]) H_exc[v]
% SOURCE    -C(1+kappa/sigma) W_2||v||^2,                    (24)
\begin{gather*}
\begin{aligned}
& E_{0,\mathrm{total}} \Vert v\Vert ^{2} \\
& +(1-C[\kappa  \sigma +\sigma ^{2}+r_{*}^{-3}]) H_{\mathrm{exc}}[v] \\
& -C(1+\kappa /\sigma ) W_{2}\Vert v\Vert ^{2},
\end{aligned}\tag{24}
\end{gather*}

with the frozen \(E_{0}\) multiplier understood pointwise in the slow
variables. This is a fast-column comparison only; the slow derivatives
in the full Gauss square have not been dropped or independently counted.

\subsubsection{6.4 Frozen spin-linear
term}\label{e-frozen-spin-linear-term}

The leading spin cross is a finite \(\sum  A_{e} p_{e}/r_{e}\), with
Y-independent Hermitian Clifford matrices
\(\Vert A_{e}\Vert \le C \kappa\). Since the \(\Omega _{e} Y_{e}\) part
of \(a_{e}=\partial _{e}+\Omega _{e} Y_{e}\) contributes a real
expectation against \(A_{e}\),

% DISPLAY e.37
% SOURCE |<v,A_e p_e v/r_e>|
% SOURCE   <= C kappa r_e^-1 sqrt(h_B,e[v]) ||v||
% SOURCE   <= epsilon h_e[v]+C kappa^2/(epsilon r_e^2)||v||^2. (25)
\begin{gather*}
\begin{aligned}
& \vert \langle v,A_{e} p_{e} v/r_{e}\rangle \vert \\
& \le  C \kappa  r_{e}^{-1} \sqrt{h_{B,e}[v]} \Vert v\Vert \\
& \le  \epsilon  h_{e}[v]+C \kappa ^{2}/(\epsilon  r_{e}^{2})\Vert v\Vert ^{2}.
\end{aligned}\tag{25}
\end{gather*}

Here \(h_{B,e}=H_{B,e}-E_{B,e}\). This also avoids any bosonic
vacuum-frequency cost. The inverse-M corrections to this frozen spin
cross pair rootwise connection errors with the bounded spin column (9),
and are bounded by \(\epsilon  H_{\mathrm{exc}}+C_{\epsilon} W_{2}\).

\subsection{7. Final scope after the
addendum}\label{e-final-scope-after-the-addendum}

Equations (2)-(16) prove the connection and Schur-loss norm statements
for arbitrary fast vectors. Equations (17)-(25) additionally prove the
signed fast-column comparison on an unrestricted oscillator core
supported in the fast tube. These estimates remove the need for a
\(\delta  \sum _{e} r_{e}\) loss at the fast-connection stage.

They still do not convert the completed, shifted slow derivative form
into a nonnegative colored internal supersymmetric form. One must
control its connection curvature/full-supercharge terms or use another
incident-weighted argument. In particular a constant loss of internal
kinetic energy cannot simply be identified with the same loss of
\(H_{a}\). That slow-form issue, the other physical potential/Yukawa
terms, and the final allocation of one non-double-counted reserve belong
to the full complementary theorem, not to this helper lemma.
